\section{Threats to validity}
\label{sec:threats}

All measurements are on one GPU (\gpuModel{}); the fitted constants are
A100-specific but the shapes should transfer. The crossover is
calibrated at a $\sim$\refPrompt-token prompt and shifts with input
length. The AR--DG quality difference ($\Delta = \qualityGap$, 95\% CI
$[-0.012, +0.011]$, $n = 300$) is equivalent at a $\pm 0.015$ margin.
This quality comparison uses the dense AR-3B baseline
(Qwen2.5-3B), not the AR-30B MoE: the 9-field F1 of
\arThreeBFnine\ was re-scored with the corrected first-object-wins parser
on the dense baseline, and the AR-30B MoE was not re-scored at the same
parser setting. The latency comparison (Table~\ref{tab:niche}) uses the
AR-30B MoE cost law, so the quality and latency arms are not the same
model.

The surviving niche has not been tested on a task whose natural output
length is $\geq 512$ tokens. The comparison in
Table~\ref{tab:niche} charges the AR arm at forced lengths for tokens that the parser
discards; the advantage in serial mode (15.6x at natural output) is a lower bound on the forced-output
comparison. Separately, the AR-30B cost law (Eq.~\ref{eq:ar}) is fitted from
batched per-event throughput, not batch-1 serial latency; the measured
batch-1 p50 at $L{=}256$ is 6148\,ms, well above the cost law's 250\,ms, so
at true batch-1 DG is faster at all measured output lengths. A task with
genuinely long natural output --- multi-record extraction, incident
reports --- would test the niche without this artifact
(Appendix~\ref{app:forced} provides an analytical projection).

The niche has not been tested against grammar-constrained AR decoding
(outlines, xgrammar, llguidance), which would eliminate over-generation,
guarantee 100\% schema validity, and remove the diffusion model's
completeness advantage. An attempt to run grammar-constrained
\ARoneb\ on vLLM 0.10.2 produced 1.0\% schema-valid output --- the
grammar enforced JSON structure but the model padded string fields it
could not answer, preventing the JSON from closing. This is a stated
limitation.

The batching ceiling reproduces across two runtimes and two model
families, but both measurements are engine-bounded
(Appendix~\ref{app:batching}). A structural argument for why the ceiling
is likely architectural --- not an engine bug --- is in
Appendix~\ref{app:batching_architectural}.

The cost-shape generalisation (``any decoder emitting fixed-size blocks
would inherit the same shape'') is inferred from one block-diffusion
model measured against one AR family. The shape claim is structural but
not tested on a second diffusion family.

DiffusionGemma-26B-A4B-it is available at
\url{https://huggingface.co/google/}\\\url{diffusiongemma-26b-a4b-it}.
All DG measurements are reproducible from this checkpoint.
